\subsection*{Graph database schema and data loading}\label{methods:schema}

GraphGWAS is built on a labelled-property graph
database\cite{robinson2015graph}; the reference implementation
uses Neo4j~5.26 Community Edition, but the schema and algorithms
are database-agnostic. The schema is
organised around nine node labels---Variant, Sample, Gene, Pathway,
GOTerm, RegulatoryElement, Population, GWASStudy,
AssociationResult---and ten edge types. \texttt{Variant} nodes carry
(variantId, chr, pos, ref, alt, af\_total, qual, gt\_packed), where
gt\_packed is a 2-bit encoded dosage array over all samples
(00~=~hom-ref, 01~=~het, 10~=~hom-alt, 11~=~missing); for
biobank-scale deployments gt\_packed is omitted and genotypes are
served from BGEN\cite{band2018bgen}. \texttt{Gene} nodes use GENCODE~v47 identifiers;
\texttt{Pathway} and \texttt{GOTerm} nodes carry name and source;
\texttt{RegulatoryElement} nodes follow ENCODE cCRE~v4;
\texttt{AssociationResult} nodes carry (run\_id, phenotype\_key,
method, beta, se, p\_value, p\_value\_log10, n\_cases, n\_controls,
maf, timestamp).

Edge types are \texttt{HAS\_CONSEQUENCE} (variant~$\rightarrow$~gene
with Variant Effect Predictor, VEP\cite{mclaren2016ensembl},
consequence), \texttt{eQTL} (variant~$\rightarrow$~gene with tissue,
slope, p-value from GTEx~v8\cite{gtex2020}), \texttt{INTERACTS\_WITH}
(gene~$\leftrightarrow$~gene with STRING combined score),
\texttt{IN\_PATHWAY}, \texttt{HAS\_GO\_TERM}, \texttt{IN\_REGULATORY},
\texttt{IN\_POPULATION}, \texttt{FOR\_VARIANT},
\texttt{IN\_STUDY}, and \texttt{NEXT}. Indexes are maintained on
(chr, pos) and variantId for Variant, symbol and geneId for Gene,
name for Pathway, and the composite (run\_id, phenotype\_key,
p\_value) for AssociationResult.

All annotation loading is performed by the
\texttt{graphgwas.annotations} CLI from publicly released input
files: GENCODE~v47 GTF\cite{frankish2023gencode}; GTEx~v8
significant eQTL tarball\cite{gtex2020}, filtered to the 49 primary
tissues; STRING~v12 \texttt{protein.physical.links} at combined
score~$\geq$~700\cite{szklarczyk2023string}; and the ENCODE cCRE~v4
registry\cite{encode2020expanded}. Every
load writes a provenance record into a \texttt{LoadManifest} node
for auditability.

\subsection*{Hybrid BGEN~+~graph~+~Pan-UKB data access}\label{methods:panukb}

All genotype-dependent methods share a single abstraction,
\texttt{load\_locus\_variants}, which returns a locus data structure
(dosage matrix, variant metadata, sample ordering) from either
(a) the graph database---unpacking \texttt{gt\_packed} for every
variant in the requested (chr, start, end) window via a
bolt-protocol query---or (b) a BGEN~v1.2 file. BGEN access uses
the \texttt{bgen} Python package directly and avoids the \texttt{.bgi}
index by pre-loading a sorted position array once per chromosome
and binary-searching locus windows. Allele frequencies computed
from BGEN dosages agree with \texttt{Variant.af\_total} in the graph
database to
$<10^{-4}$ across 2{,}364 validation variants on 1000 Genomes
chromosome~22. Conversion from VCF uses
\texttt{plink2 --export bgen-1.2 bits=8 --ref-first}. UK~Biobank
\texttt{.bgen} files are consumed by the same function without code
change.

For summary-statistics-only workloads the \texttt{graphgwas.panukb}
module streams per-locus slices of Pan-UK Biobank sumstats via
\texttt{tabix} over HTTPS against the public S3 bucket
\texttt{pan-ukb-us-east-1}; no authentication or bulk download is
required ($\sim$2.3~GB per-phenotype files accessed as
$\sim$50~kB--1~MB tabix slices per locus, $\sim$3~s each). Sumstats
(GRCh37) are lifted to GRCh38 via \texttt{pyliftover} to align with
our GENCODE~v47~/~GTEx~v8 annotation graph. The LD reference
defaults to an ancestry-matched slice of the local 1000 Genomes
BGEN; Pan-UKB in-sample LD BlockMatrices (Hail format, 47.6~TB
across six ancestries) are supported via
\texttt{panukb.fetch\_ld\_slice()} when a Hail context with
\texttt{hadoop-aws} S3 access is configured.

\subsection*{Single-locus association}

Three association tools were used across the four species,
selected by panel size, ploidy and the practical availability of
established pipelines. \textbf{Yeast (1011 Genome panel, 35 growth
traits).} Per-variant marginal effects were computed by the
internal \texttt{graphgwas.assoc} module, which fits ordinary least
squares on mean-imputed dosages with covariates (sex, ten
genomic-PCA principal components, optional population indicators)
residualised from both the phenotype $y$ and each $G_i$ before
regression. The test statistic is $z_i = \hat\beta_i / \hat s_i$
with $\hat\beta_i = (G_i^\top y - n \bar G_i \bar y) / (G_i^\top G_i - n \bar G_i^2)$
and $\hat s_i^2 = (y^\top y - \hat\beta_i (G_i^\top y - n \bar G_i \bar y)) / ((n-2)(G_i^\top G_i - n\bar G_i^2))$;
p-values derive from the $\chi^2_1$ upper tail. Population-structure
confounding is corrected by GRAMMAR+ (genomic relationship-matrix
mixed-model association with residual
re-calibration)\cite{yang2011gcta}: a genetic relatedness matrix
(GRM) built from allele frequency $> 5\%$ variants is used to
estimate heritability by restricted maximum likelihood (REML), $y$
is residualised against
GRM~$\cdot (\hat h^2 / (1 - \hat h^2))$, and per-variant statistics
are re-calibrated by $\sqrt{\lambda_{\text{GC}}}$. Across 35 yeast
traits this controls $\lambda_{\text{GC}}$ from a mean of 1.62
(range 0.85--5.04) before correction to 0.98 (range 0.90--1.02)
after, and per-variant $z$-statistics agree with PLINK2 logistic on
shared variants at $r=1.0000$ over 82{,}000 variants
(Supplementary Table~S4).

\textbf{\emph{Arabidopsis} (1001~Genomes panel, 14 priority
traits) and rice (3{,}000 Rice Genomes panel, 18 IRRI Standard
Evaluation System traits + 4 grain weight/shape traits).}
Whole-genome single-locus association was performed with
PLINK~2.0 \cite{chang2015second}. The \emph{Arabidopsis}
pipeline ran
\texttt{plink2 --vcf 1001G\_chr\{1..5\}.vcf.gz --double-id
--maf 0.01 --geno 0.05 --hwe 1e-10 --pheno phenotypes.tsv
--linear hide-covar --covar 10pcs.tsv --out
ara\_\{trait\}}, restricted to biallelic SNPs after GATK
hard-filter. The rice pipeline (3{,}024 accessions, 29.6~M
biallelic SNPs from the 3kRG pseudo-canonical VCF) ran two
multi-phenotype \texttt{plink2 --glm linear} passes on the same
genotype substrate: an 18-trait pass on the IRRI Standard
Evaluation System ordinal codes
(\texttt{tests/rice3k\_irri\_gwas\_multitrait.py}), and a
4-trait pass on the continuous grain weight (TGW) and grain
shape (GL, GW, RLW = GL/GW) phenotypes from Niu \emph{et al.}\
2021 \emph{BMC Genomics} \cite{niu2021grain} prepared by
\texttt{tests/rice3k\_grain\_shape\_phenotypes.py}
(\texttt{docs/3K.shape.phe.xlsx}: $N=2{,}453$ for GL/GW/RLW;
\texttt{docs/thousand\_grain\_weight.txt}: $N=1{,}847$ for TGW;
joint $N=1{,}832$ for any pairwise correlation analysis).
The phenotype IDs (\texttt{B}/\texttt{C}/\texttt{IRIS\_313}-prefixed
3K accession identifiers) match the plink2 \texttt{rice\_3k.psam}
directly, so no IRGC bridging is required for the grain rerun.
Phenotype QC reproduces the four-trait pairwise Pearson
correlations reported in Niu 2021 Fig.~1b to within $\Delta r <
0.04$ (e.g.\ $r(\mathrm{TGW},\mathrm{GW}){=}0.60$ vs.\ Niu's $0.63$;
$r(\mathrm{GL},\mathrm{RLW}){=}0.76$ vs.\ Niu's $0.76$;
narrow-sense $h^2$ ranges 0.88--0.93 in Niu's variance-component
decomposition). For both rice passes, all 4 (or 18) traits share
genotype I/O in a single \texttt{plink2} call. Genotypic
principal components were computed with \texttt{plink2 --pca 10}.

\textbf{Human (Pan-UK Biobank).} No GWAS was run by us.
Pre-computed Pan-UKB summary statistics from the consortium SAIGE
mixed-model pipeline\cite{karczewski2024pan,bycroft2018uk} were
consumed via \texttt{tabix}\cite{li2011tabix} over HTTPS and
lifted from GRCh37 to GRCh38 with
\texttt{pyliftover} as described in
``\nameref{methods:panukb}''.

\subsection*{Graph-augmented fine-mapping (GAFM)}\label{methods:gafm}

For each locus window ($\pm$25~kb around the lead by default; capped
at 500~kb) GAFM computes z-scores and the LD matrix $R$ from Pearson
correlations of mean-imputed dosages, then forms the LD-deconvolved
evidence
$u_i = z_i - \frac{1}{|N(i)|} \sum_{j \in N(i)} r^2_{ij} z_j$,
where $N(i) = \{j \ne i : r^2_{ij} > \tau_L\}$ with $\tau_L = 0.3$.
Negative $u_i$ are clipped to zero and the vector is rescaled so
$\max_i u_i = \max_i z_i$, preserving the statistical scale.

Functional priors are assembled in a single batched Cypher query
that returns per variant: (i) gene overlap via HAS\_CONSEQUENCE,
(ii) pathway membership, (iii) PPI partner count from
INTERACTS\_WITH at combined~score~$\geq$~700, (iv) tissue-specific
eQTL strength from GTEx~v8, and (v) conservation score. Layer
weights (1.0 per gene, 0.5 per pathway, 0.3 per PPI partner capped
at 3.0, 2.0$\times$~eqtl\_score, 1.5$\times$~max(conservation~$-$~0.5, 0))
are summed and passed through \texttt{log1p} to yield $z_{\text{func}}$,
rescaled to match $z$'s dynamic range. The combined score is
$s_i = \alpha \cdot u_i + (1-\alpha) \cdot z_{\text{func},i}$.
The default $\alpha$ is regime-specific: $\alpha{=}0.9$ for the
strong-signal six-method head-to-head benchmark
(Sup Table~S2) and the null calibration / PIP calibration
experiments; $\alpha{=}0.5$ for the weak-signal tissue-specific
eQTL regime (the 27\,--\,2 GAFM-vs-SuSiE result); and
$\alpha{=}1.0$ for the sumstats-only Pan-UK Biobank fine-mapping
path, where no eQTL functional prior is available at runtime.
A hyperparameter sensitivity sweep on the strong-signal scenario
confirms that GAFM rank-\#1 rate degrades sharply for $\alpha < 0.6$
($\alpha=0.5$: 10\%; $\alpha=0.9$: 70\%; sweep reported in Sup
Note~S2). PIPs are the softmax:
$\pi_i = \exp(s_i - \max_j s_j) / \sum_j \exp(s_j - \max_j s_j)$.
The 95\% credible set is the minimal variant set whose PIPs sum to
0.95. For the weak-signal regime, $\alpha$ is chosen by empirical
Bayes on a held-out F1 calibration set to maximise the marginal
likelihood of the observed $z$ under the mixture of statistical and
graph priors---tracking annotation informativeness locus-by-locus.
Theorem~3 (Supplementary Note~S1) shows that under linear LD decay
and bounded cross-LD between non-causal variants, $u_c > u_i$ for
all $i \ne c$; GAFM ranks the causal variant first.

\subsection*{Hierarchical belief propagation (HBP)}\label{methods:hbp}

HBP performs message passing on a three-layer factor graph
(variant, gene, pathway). The upward pass computes gene scores
$g = B_{vg}^\top b$ from the current variant belief $b$, with
$B_{vg}$ binary on HAS\_CONSEQUENCE edges (optionally weighted by
\texttt{log1p} of the eQTL composite); $g$ is then diffused one
step along the STRING PPI graph via
$g \leftarrow 0.7 \cdot g + 0.3 \cdot W_{gg} g$, where $W_{gg}$ is
the row-normalised PPI adjacency at combined~score~$\geq$~700.
Pathway scores are $p = B_{gp}^\top g$ with $B_{gp}$ binary on
IN\_PATHWAY edges. The downward pass computes
$\pi \propto B_{vg} B_{gp} p$ and $\ell_1$-normalises. The variant
belief is combined as
$b^{(t+1)} = \alpha \cdot \text{softmax}(u) + (1{-}\alpha) \cdot \pi^{(t)}$,
where $u$ is the LD-deconvolved statistic, and damped by
$b^{(t+1)} \leftarrow \lambda b^{(t)} + (1{-}\lambda) b^{(t+1)}$
before renormalisation. Defaults: $\alpha{=}0.6$, $\lambda{=}0.5$,
$T{=}5$ rounds.

\texttt{fast\_hbp\_finemap} materialises $B_{vg}$, $B_{gp}$, and
$W_{gg}$ as dense numpy arrays once per locus from a pre-loaded
graph cache; the inner loop is five matrix--vector products, giving
median runtime 0.08~s per locus on 1KG chr22 and 0.02~s on smaller
windows. Theorem~2 shows the HBP update is a strict $\ell_1$
contraction with rate
$L = \lambda + (1-\lambda)(1-\alpha) \cdot \rho(B_{vg} B_{gp} B_{gp}^\top B_{vg}^\top) < 1$
whenever $\lambda > 0$, so by Banach's fixed-point theorem HBP
converges geometrically to a unique self-consistent fixed point.

\subsection*{Baseline methods}

SuSiE\cite{wang2020simple,zou2022fine} is run through R
\texttt{susieR}~v0.14.2 via \texttt{susie\_suff\_stat} with
$L{=}10$ single effects and
\texttt{estimate\_residual\_variance\,=\,TRUE}, on the same $R$ and
$z$ inputs as GAFM and HBP.
FINEMAP\cite{benner2016finemap}~v1.4.2 is invoked as a subprocess
with precomputed \texttt{.z} and \texttt{.ld} files and flags
\texttt{-{-}sss -{-}n-causal-snps 5 -{-}n-iter 100000}; per-variant
PIPs are parsed from \texttt{.snp} output. SuSiE-inf and FINEMAP-inf
use the FinucaneLab Python implementations
(\texttt{susieinf}~v1.4, \texttt{finemapinf}~v1.3)\cite{cui2024improving}
with default infinitesimal-variance hyperparameters. The
Polyfun-proxy\cite{weissbrod2020polyfun} sets per-variant prior
weights to $w_i = \log(1 + n_{\text{eqtl\_edges},i})$ normalised per
locus, passed as \texttt{prior\_weights}.
SBayesRC\cite{zheng2024leveraging}~v0.2.6 was compiled with
CXX17 and \texttt{BOOST\_ALLOW\_DEPRECATED\_HEADERS} in
\texttt{src/Makevars}, using the published EUR HapMap~3 LD
reference (3~GB) and Baseline~2.2 stratified LD score regression
annotations (1.9~GB)\cite{finucane2015partitioning,buliksullivan2015ld};
Markov chain Monte Carlo (MCMC) ran for 500 iterations
(250 burn-in). All six baselines are wrapped in a
common Python interface that consumes the same locus data structure
as GAFM and HBP.

\subsection*{Simulation protocols}

Simulations are driven by \texttt{graphgwas.simulate}. Phenotypes
are generated as $y = G\beta + \epsilon$, $\epsilon \sim
\mathcal{N}(0, \sigma^2 I)$ scaled so $\text{var}(G\beta)/\text{var}(y)$
equals the target $h^2$. Each replicate logs its seed, causal-variant
identity, and sampled window.

\textbf{F1 (single causal in LD).} Random 50~kb windows on chr22,
causal drawn from AF~$\in$~[0.05, 0.50], $h^2 \in \{0.02, 0.05, 0.10, 0.20\}$,
30--200 replicates per cell.
\textbf{S1 (pure interaction).} Two causal variants placed
$>$~1~Mb apart so $r^2 \approx 0$, $\beta_1 = \beta_2 = 0$,
$\beta_I = 1.5$, $h^2 = 0.30$.
\textbf{F1-eQTL.} Causal variants restricted to GTEx~v8 significant
eQTLs with tissue-specificity $-\log_{10} p > 15$ in $\leq$~2
tissues; $\beta{=}0.15$, $h^2{=}0.01$; 79 replicates provide the
weak-signal headline.
\textbf{100-rep weak-signal replication.} 30 rotating centres across
chr22 17--46~Mb, $\beta{=}0.2$, $h^2{=}0.02$.
\textbf{Null.} $\beta{=}0$ everywhere, Gaussian phenotype, 100
replicates on 50~kb yeast windows.
\textbf{Cross-ancestry.} 1KG restricted to EUR ($n{=}503$),
AFR~($n{=}661$), EAS~($n{=}504$); F1, $h^2{=}0.05$, 30 replicates
per ancestry.
\textbf{Power-vs-$N$.} 1KG subsampled at
$N \in \{500, 1{,}000, 2{,}000, 3{,}000\}$, F1 with $h^2{=}0.10$.
\textbf{Cross-species.} Real phenotypes on \emph{Arabidopsis}
1001G FT10 (1{,}003 accessions), yeast 1011G 35 growth traits
(971 strains after matching), 18 IRRI Standard Evaluation System
trait codes from the 3{,}000 Rice Genomes\cite{wang2018genomic}
phenotype catalogue (2{,}266 accessions linked to plink2 sample
IDs through the IRGC accession cross-reference), and four Pan-UKB
phenotypes (BMI, height, LDL, triglycerides; sumstats fetched via
\texttt{tabix} over HTTPS, GRCh37$\to$GRCh38 lifted with
\texttt{pyliftover}).

\subsection*{Cross-species fine-mapping protocol}\label{methods:multispecies}

The same end-to-end pipeline was applied to all four species, with
only the loaded annotation graph differing:
\textbf{rice} --- RAP-DB / MSU gene
model\cite{sakai2013rapdb} from
snpEff\cite{cingolani2012snpeff} ANN tags on the 3kRG
pseudo-canonical VCF + Ren 2023\cite{ren2023rice} 269-gene
grain-quality pathway labels + RicePPINet\cite{liu2017riceppinet}
PPI at probability $\geq 0.7$;
\textbf{yeast} --- SGD\cite{engel2014sgd} ORFs and
GO\cite{ashburner2000go}-slim annotations from the SGD features
catalogue, with BIOGRID\cite{oughtred2021biogrid} 5.0.256
\emph{S.~cerevisiae} S288c physical interactions (filtered to
263{,}403 high-confidence edges across 5{,}813 ORFs);
\textbf{Arabidopsis} --- TAIR10\cite{berardini2015tair} gene model
from NCBI, Plant Reactome\cite{gillespie2022reactome} pathway
membership, and STRING-\emph{Ath} v12\cite{szklarczyk2023string}
PPI at score $\geq 700$ resolved through
UniProt\cite{uniprot2023} \texttt{ARATH\_3702} idmapping
(\texttt{Gene\_OrderedLocusName} field) to recover
20{,}192 \emph{Arabidopsis} genes with at least
one high-confidence PPI partner;
\textbf{human} --- per-chromosome GTEx v8\cite{gtex2020} eQTL
across 49 tissues (4.59~M annotated variants), STRING
v12\cite{szklarczyk2023string} \emph{Homo sapiens} PPI at score
$\geq 700$, and GENCODE v47\cite{frankish2023gencode} ENSG--symbol
resolution.

For each species, lead loci were called by greedy clustering of
genome-wide-significant variants ($p \leq 5\times10^{-8}$) at
trait-LD-appropriate windows (15~kb yeast; 100~kb \emph{Arabidopsis};
250~kb rice IRRI / rice grain;
500~kb human). Up to four leads per trait were
fine-mapped within $\pm$15--500~kb LD windows using the
species-specific 1KG-equivalent reference panel for LD computation
(yeast 1011G; \emph{Arabidopsis} 1001G; rice 3kRG;
1KG NYGC GRCh38 for human). For the rice grain weight + shape pass, a 5-method panel was
applied at the top 5 genome-wide-significant + top 2 suggestive
leads per trait (28 loci total) within $\pm$100\,kb LD windows
after filtering to
common variants (MAF $\geq 0.05$, matching Niu \emph{et al.}\ 2021's
fine-mapping convention): GAFM and HBP from this work; SuSiE
\cite{wang2020simple} via \texttt{susieR::susie\_rss} on signed
in-sample LD; SuSiE-inf and FINEMAP-inf
\cite{cui2024improving} as the FinucaneLab Python reference
implementations of the two state-of-the-art infinitesimal-effects
fine-mappers from Wu \emph{et al.}\ 2026 (Fig.~4b). All five
methods used the same per-locus z-statistics (BETA/SE), the same
in-sample 3kRG LD matrix, and the same 95\% credible-set
coverage; GAFM and HBP additionally use the rice multi-omics graph
prior (RAP-DB variant--gene + Ren 2023 gene--pathway +
RicePPINet PPI). SBayesRC \cite{zheng2024leveraging} is included as the 6th method.
The public SBayesRC LD reference is built from human HapMap3
European samples and is therefore incompatible with the rice
3kRG cohort (different species, different LD structure). We
constructed a 3kRG-specific LD reference from scratch using
gctb 2.05beta and the R
SBayesRC \texttt{LDstep1}--\texttt{LDstep4} pipeline. We define
23 non-overlapping blocks covering all 41 fine-mapped leads
(merging $\pm$250\,kb intervals around each lead; mean block
size $\sim$600\,kb, range 500\,kb--1.1\,Mb). Block construction
restricts to the 188\,K MAF~$\geq$~0.05 variants in those windows
extracted from the 3kRG \texttt{.pgen}; gctb computes a per-block
full LD matrix (\texttt{--make-full-ldm}); SBayesRC's
\texttt{LDstep3} eigen-decomposes each block at variance
threshold 0.995. SBayesRC is then run per trait at default
\texttt{niter=3000, burn=1000} with the four-component mixture
prior (\texttt{startPi=(0.99, 0.005, 0.003, 0.001, 0.001)},
\texttt{gamma=(0, 0.001, 0.01, 0.1, 1)}). Per-variant posterior
PIPs are extracted at the 41 fine-map windows and merged into
the per-locus 6-method scorecard.

\textbf{Mixture-prior posterior reweighting and ensemble (GAFM-MX, HBP-MX, ENS).}
The 5-method panel above revealed that GAFM and HBP, while
200--700$\times$ faster than SuSiE on this rice grain pass, gave large
credible sets (median CS $\approx$ 1{,}500) on secondary loci under the
strong inflation regime ($\lambda_{GC}{=}1.41$--$1.78$ from XI/GJ
structure under PC1--PC10). SBayesRC's continuous-shrinkage mixture
prior dominated those secondary loci, suggesting that the same
effect-size-distribution correction would help GAFM/HBP. We therefore
add three corrections, all applied by default to GAFM-MX, HBP-MX and
ENS:
\emph{(i)} \textbf{$\lambda_{GC}$ deflation} of z-scores at the trait
level ($z \to z / \sqrt{\lambda_{GC}}$, with the matching SE inflation in
the SBayesRC \texttt{.ma} files); this is the standard genomic-control
correction and is also applied to all six original methods.
\emph{(ii)} \textbf{Mixture-prior posterior reweighting (MX)}: the
per-variant PIPs from GAFM and HBP are multiplied by a 4-component
Wakefield mixture-BF on the
\emph{LD-deconvolved}, $\lambda_{GC}$-deflated z-scores
($\pi=(0.005, 0.003, 0.001, 0.001)$ renormalised within the non-zero
components; $\gamma=(0.001, 0.01, 0.1, 1.0)$; $V_k = N\gamma_k$;
ABF$_k(z) = (1+V_k)^{-1/2} \exp(z^2 V_k / (2(1+V_k)))$). LD
deconvolution (the same step GAFM and HBP use to build their base
PIPs; Methods \S\,``LD-aware unique-statistics deconvolution'')
removes neighbour-driven inflation in $z$ before the mixture so
that LD-correlated noise variants with marginally larger marginal
$|z|$ do not outcompete the causal during multiplicative
reweighting. The reweighted distribution preserves
$\sum_i \mathrm{PIP}_i$ (expected number of causals unchanged) but
redistributes mass toward large-$|z|$ variants. The two new
methods are written GAFM-MX and HBP-MX.
\emph{(iii)} \textbf{Ensemble (ENS)}: the mean of GAFM-MX and HBP-MX
PIPs per variant, with credible sets re-derived from the averaged PIPs
at 95\% coverage. The combined panel is therefore nine methods:
GAFM, HBP, GAFM-MX, HBP-MX, ENS, SuSiE, SuSiE-inf, FINEMAP-inf
and SBayesRC. The full nine-method per-locus credible-set
comparison is given in Supplementary Table~S9.
Validation used 250~kb-window overlap
against literature catalogues: Bloom 2015\cite{bloom2015genetic} +
Peter 2018\cite{peter2018genome} for yeast,
AraGWAS\cite{togninalli2018aragwas} Bonferroni hits + canonical
flowering-time genes for \emph{Arabidopsis},
Ren 2023\cite{ren2023rice} for rice, and the
GWAS-Catalog\cite{buniello2019nhgri} plus
lipid-genetics literature for human chromosome 22 (the genome-wide
human catalogue used here lists 36 loci across BMI, height, LDL
direct and triglycerides, drawn from GIANT 2018, GLGC 2013, Klarin
2018 and Yengo 2022). To calibrate the validation rate against
chance, we ran a chromosome-matched null-overlap permutation: for
each species, draw 1{,}000 random ``leads'' matched to the
chromosome distribution of the observed lead set and compute the
same 250\,kb-window overlap rate against the catalogue. The
empirical p-value of the observed rate is reported in Sup Note~S2.
Rice (33.3\% observed vs 0.9\% null mean; $p<0.001$; 38.4-fold
enrichment) and \emph{Arabidopsis} (5.6\% observed vs 0.0\% null
mean; $p<0.001$; catalogue too sparse for an enrichment estimate)
are significantly above chance; yeast and human require
coordinate-level catalogue files that are not yet integrated in
the present analysis.

\subsection*{Heterogeneity-by-intervention experiment}\label{methods:intervention}

The graph cache schema was extended with an optional per-variant
continuous \texttt{prior\_score} field. HBP re-mixes
\texttt{prior\_score} into the final posterior after belief
propagation as $\textsf{PIP} = 0.5 \cdot \textsf{PIP}_{\text{BP}} +
0.5 \cdot \mathrm{softmax}(\alpha \cdot u + (1-\alpha) \cdot
(z_{\text{func}} + 0.5 \cdot s))$ with $s$ the per-variant
\texttt{prior\_score} rescaled to z-score range and $u$ the
LD-deconvolved unique signal. GAFM adds the rescaled
\texttt{prior\_score} directly to its $z_{\text{func}}$ mixture
before LD deconvolution. Caches without
\texttt{prior\_score} reproduce baseline results bit-for-bit.

The smallest readily available heterogeneous layer was added per
species: for human, ENCODE cCRE class flags (926{,}535 elements
across 9 classes; PLS-CTCF $=$ 1.5, PLS $=$ 1.0, pELS-CTCF $=$ 1.2,
pELS $=$ 0.8, dELS-CTCF $=$ 1.0, dELS $=$ 0.6, CTCF-only $=$ 0.7,
DNase-H3K4me3-CTCF $=$ 0.9, DNase-H3K4me3 $=$ 0.5);
for yeast, ORF-feature-class scores (verified ORF $=$ 1.0,
tRNA/ncRNA $=$ 0.7, pseudogene $=$ 0.4) plus the BIOGRID PPI layer;
for \emph{Arabidopsis}, TAIR10 feature-class scores
(CDS $=$ 1.0, exon $=$ 0.8, gene with $\pm$2-kb promoter $=$ 0.5);
for rice, snpEff impact-class scores from the pseudo-canonical
VCF re-scan (HIGH $=$ 1.0, MODERATE $=$ 0.7, LOW $=$ 0.4,
MODIFIER $=$ 0.1). The same fine-mapping pipeline as above was
re-run on identical lead-loci sets; the only change between
baseline and intervention runs is the loaded cache.

\subsection*{Evaluation and hardware}

For each replicate we record the rank of the causal variant by PIP,
the PIP value, the 95\% credible-set size, and wall-clock runtime.
Summary statistics are rank-\#1 rate, mean rank, mean PIP,
head-to-head win counts (ties broken by rank), and the sign-test
$p$-value on asymmetric head-to-head outcomes. PIP calibration is
reported as observed true-discovery rate in eight bins of width
0.125 across the pooled 200-rep~$\times$~4-$h^2$ design. Null FPR
is the fraction of replicates with $\max$\,PIP~$>$~0.5. All runtimes
were measured on a single workstation---Intel Core i9-13900K
(24~cores), 64~GB DDR5, Ubuntu 24.04, Python 3.13, graph database
heap 16~GB / page cache 16~GB---and report median wall-clock per locus
excluding graph-bootstrap and LD-reference load.

\subsection*{Software and reproducibility}

GraphGWAS source is released under the MIT licence at
\url{https://github.com/jfmao/GraphGWAS}. Python dependencies:
\texttt{numpy}\cite{harris2020numpy}, \texttt{scipy}\cite{virtanen2020scipy},
\texttt{pandas}, \texttt{bgen}~1.9.9, \texttt{bgen-reader}~4.0.9,
\texttt{neo4j}~6.1, \texttt{susieinf}~1.4, \texttt{finemapinf}~1.3,
\texttt{pyliftover}~0.4. R: \texttt{susieR}~0.14.2,
\texttt{SBayesRC}~0.2.6. External binaries:
\texttt{plink2}\cite{chang2015second}~v2.0.0-a.6.5LM, FINEMAP~v1.4.2,
\texttt{bcftools}\cite{danecek2021twelve}~1.19,
\texttt{tabix}\cite{li2011tabix}~(htslib~1.23.1), GCTB~2.5. Optional
Hail\cite{hailteam2024}~0.2.138 in a dedicated Python~3.11 env for
Pan-UKB LD BlockMatrix access. The Model Context Protocol (MCP)
server\cite{mcp2024} is implemented with the official Python SDK;
the graph neural network layer uses PyTorch
Geometric\cite{fey2019fast}.
Benchmark JSONs, figure scripts, and pre-built graph-database
dumps (Neo4j~5.26 format; yeast 0.5~GB, human with multi-omics
17~GB) are archived alongside the code. Every
figure and table is regeneratable by a single command documented in
\texttt{docs/REPRODUCIBILITY.md}.

The fine-mapping methods rigorously benchmarked in this paper are
HBP and GAFM; CLGF (cross-locus EM) and GLEM (graph-latent-embedding
fine-mapping) are implemented in the codebase but not benchmarked
here. Together they are one method class within the broader
GraphGWAS platform.
A full description of the platform---including methods implemented
for heritability estimation, multivariate and cross-trait analysis,
polygenic risk scoring, Mendelian randomisation and
gene--environment interaction, none of which are benchmarked in
the present paper---together with an honest benchmark-status table
for each method class, is provided in Supplementary Note~S2. The
44-command CLI hierarchy is visualised as Supplementary
Figure~S3; a per-command reference manual is shipped at
\texttt{docs/manual/index.md}, and an end-to-end vignette
reproducing the FTO/BMI headline result from Pan-UKB summary
statistics in $\sim$15 minutes is provided at
\texttt{vignettes/fine-mapping-quickstart.md}. The same procedures
are exposed through a 37-endpoint FastAPI REST server
(\texttt{graphgwas serve}) and a 16-tool MCP server
(\texttt{graphgwas mcp}) for programmatic and AI-agent access.
